\subsection{Independent audit of the Deligne bound comparison}\label{endpoint-purity_review--independent-audit-of-the-deligne-bound-comparison}

Date: 13 September 2026. This is a complete written review of the operator, coefficient, period, purity, and critical-kernel parts of the delivered \texttt{Tau\_Deligne\_Bound\_Audit/NOTE.tex}. The complete note was read, including its Gamma argument; independent verification of the quantitative Gamma estimate belongs to the separate Gamma review. This document supplies full proofs for its new comparison statements. It preserves the delivered note unchanged.

\subsubsection{Finding and precise scope}\label{endpoint-purity_review--finding-and-precise-scope}

The delivered note proves a failure of a particular implication from auxiliary exponential Frobenius purity to the original arithmetic operator estimate. Its rank-one and reflection-stable counterexamples survive the review below. They do not prove failure of the Split-Zero or broader F1 geometry program. Their test polynomials are not known to be factors of the actual function ENDPOINTSOURCEF8C6D3817D775A0F5921DFECSLOT0END, and the exact arithmetic comparison records that difference by multiplication by ENDPOINTSOURCEF8C6D3817D775A0F5921DFECSLOT1END.

Two narrow corrections are required in the exposition. The rank and trace statements attached to the Artin--Schreier equation must specify the nontrivial additive-character sheaf, or equivalently its character summand in the cover. The whole affine cover also has a degree-two compact-support contribution. Furthermore, vanishing of ENDPOINTSOURCEF8C6D3817D775A0F5921DFECSLOT2END and entire divisibility by ENDPOINTSOURCEF8C6D3817D775A0F5921DFECSLOT3END do not require that every selected zero order be its full actual order; the extra full-order condition makes ENDPOINTSOURCEF8C6D3817D775A0F5921DFECSLOT4END a unit. Neither correction invalidates the exhibited operator counterexamples.

The strongest comparisons proved here include a complete classification of the possible pure-target intertwiners, with nilpotent lengths and exponential collisions retained; a universal observable quotient for each specified target Frobenius; and a tensor construction from the original arithmetic packet whose finite-field realization has the correct tensor weight ENDPOINTSOURCEF8C6D3817D775A0F5921DFECSLOT5END. The latter has an explicit injective complex period map retaining the original Taylor unit and sum action. Its existence prevents a failure of the simpler rank-ENDPOINTSOURCEF8C6D3817D775A0F5921DFECSLOT6END, weight-one auxiliary inference from being mistaken for failure of every geometric comparison.

\subsubsection{1. The actual coefficient complex and its recovered operator}\label{endpoint-purity_review--the-actual-coefficient-complex-and-its-recovered-operator}

Retain the original monic annihilator

ENDPOINTSOURCEF8C6D3817D775A0F5921DFECSLOT7END

and the specified potential

ENDPOINTSOURCEF8C6D3817D775A0F5921DFECSLOT8END

Let ENDPOINTSOURCEF8C6D3817D775A0F5921DFECSLOT9END. The two-term complex in degrees zero and one is

ENDPOINTSOURCEF8C6D3817D775A0F5921DFECSLOT10END

For a nonzero polynomial ENDPOINTSOURCEF8C6D3817D775A0F5921DFECSLOT11END of degree ENDPOINTSOURCEF8C6D3817D775A0F5921DFECSLOT12END and leading coefficient ENDPOINTSOURCEF8C6D3817D775A0F5921DFECSLOT13END, the term of degree ENDPOINTSOURCEF8C6D3817D775A0F5921DFECSLOT14END in ENDPOINTSOURCEF8C6D3817D775A0F5921DFECSLOT15END is exactly ENDPOINTSOURCEF8C6D3817D775A0F5921DFECSLOT16END. Every other term has smaller degree. Thus ENDPOINTSOURCEF8C6D3817D775A0F5921DFECSLOT17END is injective. For a polynomial ENDPOINTSOURCEF8C6D3817D775A0F5921DFECSLOT18END of degree at least ENDPOINTSOURCEF8C6D3817D775A0F5921DFECSLOT19END, subtract the image under ENDPOINTSOURCEF8C6D3817D775A0F5921DFECSLOT20END of its leading coefficient times ENDPOINTSOURCEF8C6D3817D775A0F5921DFECSLOT21END. The degree strictly decreases. The procedure terminates and writes ENDPOINTSOURCEF8C6D3817D775A0F5921DFECSLOT22END with ENDPOINTSOURCEF8C6D3817D775A0F5921DFECSLOT23END. Uniqueness follows because a nonzero image has degree at least ENDPOINTSOURCEF8C6D3817D775A0F5921DFECSLOT24END. Consequently

ENDPOINTSOURCEF8C6D3817D775A0F5921DFECSLOT25END

This proves freeness over the displayed coefficient ring and its compatibility with every coefficient specialization. The special fibre at ENDPOINTSOURCEF8C6D3817D775A0F5921DFECSLOT26END is exactly ENDPOINTSOURCEF8C6D3817D775A0F5921DFECSLOT27END. No repeated root is removed.

On ENDPOINTSOURCEF8C6D3817D775A0F5921DFECSLOT28END, put ENDPOINTSOURCEF8C6D3817D775A0F5921DFECSLOT29END on polynomial coefficients in both terms. Directly,

ENDPOINTSOURCEF8C6D3817D775A0F5921DFECSLOT30END

Their sum is zero, so ENDPOINTSOURCEF8C6D3817D775A0F5921DFECSLOT31END descends to ENDPOINTSOURCEF8C6D3817D775A0F5921DFECSLOT32END. In the fixed increasing-power frame, reduction of ENDPOINTSOURCEF8C6D3817D775A0F5921DFECSLOT33END uses only ENDPOINTSOURCEF8C6D3817D775A0F5921DFECSLOT34END; therefore the matrix is

ENDPOINTSOURCEF8C6D3817D775A0F5921DFECSLOT35END

The operator ENDPOINTSOURCEF8C6D3817D775A0F5921DFECSLOT36END preserves ENDPOINTSOURCEF8C6D3817D775A0F5921DFECSLOT37END: its derivative on ENDPOINTSOURCEF8C6D3817D775A0F5921DFECSLOT38END contributes ENDPOINTSOURCEF8C6D3817D775A0F5921DFECSLOT39END, which belongs to that ideal. Hence it induces precisely ENDPOINTSOURCEF8C6D3817D775A0F5921DFECSLOT40END on the special fibre. This is a regular operator specialization even though ENDPOINTSOURCEF8C6D3817D775A0F5921DFECSLOT41END itself has a pole. Multiplication by ENDPOINTSOURCEF8C6D3817D775A0F5921DFECSLOT42END on arbitrary representatives does not separately define an algebra structure on ENDPOINTSOURCEF8C6D3817D775A0F5921DFECSLOT43END; the displayed connection and fixed-frame endomorphism give its exact replacement.

\subsubsection{2. Rank-one periods, orientation, and the finite-field trace}\label{endpoint-purity_review--rank-one-periods-orientation-and-the-finite-field-trace}

For ENDPOINTSOURCEF8C6D3817D775A0F5921DFECSLOT44END, with ENDPOINTSOURCEF8C6D3817D775A0F5921DFECSLOT45END, the coefficient ring ENDPOINTSOURCEF8C6D3817D775A0F5921DFECSLOT46END contains every coefficient. Let ENDPOINTSOURCEF8C6D3817D775A0F5921DFECSLOT47END, put ENDPOINTSOURCEF8C6D3817D775A0F5921DFECSLOT48END, and first take ENDPOINTSOURCEF8C6D3817D775A0F5921DFECSLOT49END real. On the upward contour ENDPOINTSOURCEF8C6D3817D775A0F5921DFECSLOT50END, ENDPOINTSOURCEF8C6D3817D775A0F5921DFECSLOT51END,

ENDPOINTSOURCEF8C6D3817D775A0F5921DFECSLOT52END

Write the real integral as ENDPOINTSOURCEF8C6D3817D775A0F5921DFECSLOT53END. Differentiation is justified by a Gaussian bound. Integration by parts, whose endpoint terms vanish, gives

ENDPOINTSOURCEF8C6D3817D775A0F5921DFECSLOT54END

Thus ENDPOINTSOURCEF8C6D3817D775A0F5921DFECSLOT55END. The Gaussian value ENDPOINTSOURCEF8C6D3817D775A0F5921DFECSLOT56END, obtained by squaring the integral and using polar coordinates, yields

ENDPOINTSOURCEF8C6D3817D775A0F5921DFECSLOT57END

Both sides of the unsquared period identity are entire in ENDPOINTSOURCEF8C6D3817D775A0F5921DFECSLOT58END: compact sets of complex ENDPOINTSOURCEF8C6D3817D775A0F5921DFECSLOT59END still admit an integrable Gaussian majorant. Hence the unsquared identity extends to complex ENDPOINTSOURCEF8C6D3817D775A0F5921DFECSLOT60END; the last real norm formula is asserted only for real ENDPOINTSOURCEF8C6D3817D775A0F5921DFECSLOT61END.

For the inherited ray order, ENDPOINTSOURCEF8C6D3817D775A0F5921DFECSLOT62END, ENDPOINTSOURCEF8C6D3817D775A0F5921DFECSLOT63END, ENDPOINTSOURCEF8C6D3817D775A0F5921DFECSLOT64END, and ENDPOINTSOURCEF8C6D3817D775A0F5921DFECSLOT65END. Its path travels from ENDPOINTSOURCEF8C6D3817D775A0F5921DFECSLOT66END to zero and then to ENDPOINTSOURCEF8C6D3817D775A0F5921DFECSLOT67END. Thus ENDPOINTSOURCEF8C6D3817D775A0F5921DFECSLOT68END as oriented chains and ENDPOINTSOURCEF8C6D3817D775A0F5921DFECSLOT69END. The unsquared sign in the delivered note is correct.

Now fix a finite field ENDPOINTSOURCEF8C6D3817D775A0F5921DFECSLOT70END of odd characteristic, a nontrivial complex additive character ENDPOINTSOURCEF8C6D3817D775A0F5921DFECSLOT71END, ENDPOINTSOURCEF8C6D3817D775A0F5921DFECSLOT72END, and ENDPOINTSOURCEF8C6D3817D775A0F5921DFECSLOT73END. Translation ENDPOINTSOURCEF8C6D3817D775A0F5921DFECSLOT74END is a bijection and gives

ENDPOINTSOURCEF8C6D3817D775A0F5921DFECSLOT75END

The character takes values in roots of unity, so its conjugate at ENDPOINTSOURCEF8C6D3817D775A0F5921DFECSLOT76END is its value at ENDPOINTSOURCEF8C6D3817D775A0F5921DFECSLOT77END. The invertible coordinate change

ENDPOINTSOURCEF8C6D3817D775A0F5921DFECSLOT78END

therefore gives

ENDPOINTSOURCEF8C6D3817D775A0F5921DFECSLOT79END

Indeed the inner sum is ENDPOINTSOURCEF8C6D3817D775A0F5921DFECSLOT80END at ENDPOINTSOURCEF8C6D3817D775A0F5921DFECSLOT81END. For ENDPOINTSOURCEF8C6D3817D775A0F5921DFECSLOT82END, choose ENDPOINTSOURCEF8C6D3817D775A0F5921DFECSLOT83END with ENDPOINTSOURCEF8C6D3817D775A0F5921DFECSLOT84END; shifting ENDPOINTSOURCEF8C6D3817D775A0F5921DFECSLOT85END by ENDPOINTSOURCEF8C6D3817D775A0F5921DFECSLOT86END multiplies that inner sum by a number unequal to one, so the inner sum vanishes.

The cohomological object for this sum is

ENDPOINTSOURCEF8C6D3817D775A0F5921DFECSLOT87END

With the geometric-Frobenius convention in the source, its trace is the negative of DP.7. The degree-two character-polynomial theorem makes DP.9 one-dimensional and all other compact cohomology of this \textbf{character sheaf} zero. Thus the sole Frobenius eigenvalue has modulus ENDPOINTSOURCEF8C6D3817D775A0F5921DFECSLOT88END, including its trace sign.

For precision about the cover, fix ENDPOINTSOURCEF8C6D3817D775A0F5921DFECSLOT89END and write a nontrivial ENDPOINTSOURCEF8C6D3817D775A0F5921DFECSLOT90END as ENDPOINTSOURCEF8C6D3817D775A0F5921DFECSLOT91END, ENDPOINTSOURCEF8C6D3817D775A0F5921DFECSLOT92END. The cover used for this summand is ENDPOINTSOURCEF8C6D3817D775A0F5921DFECSLOT93END; alternatively fix ENDPOINTSOURCEF8C6D3817D775A0F5921DFECSLOT94END and the corresponding character. The finite étale pushforward from that cover decomposes by the character idempotents of its translation group ENDPOINTSOURCEF8C6D3817D775A0F5921DFECSLOT95END. Its trivial summand is the constant sheaf on ENDPOINTSOURCEF8C6D3817D775A0F5921DFECSLOT96END, whose ENDPOINTSOURCEF8C6D3817D775A0F5921DFECSLOT97END is ENDPOINTSOURCEF8C6D3817D775A0F5921DFECSLOT98END. Every nontrivial character summand has the type DP.9. Consequently the whole cover has an additional degree-two contribution, whereas DP.9 has the rank and alternating-trace sign used in the note. This is the required clarification of its equations (14) and (22)--(23).

Finally, on the recovered complex one-dimensional source ENDPOINTSOURCEF8C6D3817D775A0F5921DFECSLOT99END, any positive scalar Gram ENDPOINTSOURCEF8C6D3817D775A0F5921DFECSLOT100END gives

ENDPOINTSOURCEF8C6D3817D775A0F5921DFECSLOT101END

This identity and DP.8 coexist in the same coefficient family. Hence the indicated coefficient specialization and period identity do not imply a vanishing arithmetic weight-one defect.

\subsubsection{3. Primary-source hypotheses and the reflection-stable example}\label{endpoint-purity_review--primary-source-hypotheses-and-the-reflection-stable-example}

The local original Deligne scan was inspected visually on complete printed pages 215 and 216, including the theorem, its character-sheaf definition, and its family argument. The complete relevant S20 French text and the dated primary correction were also read. The primary bibliographic page was checked online: \href{https://www.numdam.org/item/PMIHES_1980__52__137_0/}{Deligne, La conjecture de Weil II}.

The theorem used is Weil II (3.7.2.3): for a degree-ENDPOINTSOURCEF8C6D3817D775A0F5921DFECSLOT102END polynomial in ENDPOINTSOURCEF8C6D3817D775A0F5921DFECSLOT103END variables, ENDPOINTSOURCEF8C6D3817D775A0F5921DFECSLOT104END, and smooth leading projective hypersurface, the associated nontrivial additive-character sheaf has compact cohomology concentrated in degree ENDPOINTSOURCEF8C6D3817D775A0F5921DFECSLOT105END, of dimension ENDPOINTSOURCEF8C6D3817D775A0F5921DFECSLOT106END, pure of weight ENDPOINTSOURCEF8C6D3817D775A0F5921DFECSLOT107END. The parameter sheaves are lisse by (3.7.3). Its propagation argument uses mixedness as well as lissity in (1.8.12). In (3.7.4) the final reference is Roman \textbf{I.8.11}. This audit does not claim a fresh proof of Deligne's theorem or a visual rereading of all Weil II pages.

For DP.1, ENDPOINTSOURCEF8C6D3817D775A0F5921DFECSLOT108END, ENDPOINTSOURCEF8C6D3817D775A0F5921DFECSLOT109END, and ENDPOINTSOURCEF8C6D3817D775A0F5921DFECSLOT110END, the leading coefficient is exactly ENDPOINTSOURCEF8C6D3817D775A0F5921DFECSLOT111END. In ENDPOINTSOURCEF8C6D3817D775A0F5921DFECSLOT112END, its zero scheme is empty because that coefficient is invertible. The empty scheme is smooth. Taking ENDPOINTSOURCEF8C6D3817D775A0F5921DFECSLOT113END is stronger than the required ENDPOINTSOURCEF8C6D3817D775A0F5921DFECSLOT114END and retains all denominators. The lower critical points may collide; they do not alter this leading homogeneous hypothesis. Thus the stated theorem applies after each specified good coefficient reduction.

The proposed quartet polynomial is correct. With ENDPOINTSOURCEF8C6D3817D775A0F5921DFECSLOT115END, its four factors have roots ENDPOINTSOURCEF8C6D3817D775A0F5921DFECSLOT116END. Pairing conjugates at each real part gives

ENDPOINTSOURCEF8C6D3817D775A0F5921DFECSLOT117END

It is monic over ENDPOINTSOURCEF8C6D3817D775A0F5921DFECSLOT118END, real, and invariant under ENDPOINTSOURCEF8C6D3817D775A0F5921DFECSLOT119END. Its roots have real parts ENDPOINTSOURCEF8C6D3817D775A0F5921DFECSLOT120END, including both signs of their nonzero imaginary parts. Its degree-five potential has leading coefficient ENDPOINTSOURCEF8C6D3817D775A0F5921DFECSLOT121END. For all specified good reductions with ENDPOINTSOURCEF8C6D3817D775A0F5921DFECSLOT122END, the previous theorem applies to its character sheaf and yields weight one. Thus reflection symmetry does not restore the failed operator implication. The statement does not assert that ENDPOINTSOURCEF8C6D3817D775A0F5921DFECSLOT123END divides ENDPOINTSOURCEF8C6D3817D775A0F5921DFECSLOT124END.

The exact translation proof in the retained source was also read in full. For ENDPOINTSOURCEF8C6D3817D775A0F5921DFECSLOT125END, ENDPOINTSOURCEF8C6D3817D775A0F5921DFECSLOT126END, and ENDPOINTSOURCEF8C6D3817D775A0F5921DFECSLOT127END, direct substitution gives the chain isomorphism and

ENDPOINTSOURCEF8C6D3817D775A0F5921DFECSLOT128END

The fixed-frame matrix ENDPOINTSOURCEF8C6D3817D775A0F5921DFECSLOT129END for ENDPOINTSOURCEF8C6D3817D775A0F5921DFECSLOT130END has inverse ENDPOINTSOURCEF8C6D3817D775A0F5921DFECSLOT131END. The scalar is nonvanishing on ENDPOINTSOURCEF8C6D3817D775A0F5921DFECSLOT132END and is retained. For a holomorphic determinant-line metric multiplied by ENDPOINTSOURCEF8C6D3817D775A0F5921DFECSLOT133END, the curvature change is ENDPOINTSOURCEF8C6D3817D775A0F5921DFECSLOT134END locally; logarithm branches differ by constants and give the same derivative. The corresponding additive trace defect changes by ENDPOINTSOURCEF8C6D3817D775A0F5921DFECSLOT135END. The reflection-preserving DP.11 shows that the broader counterexample does not rely on changing the original reflection centre by DP.12.

\subsubsection{4. Classification of all intertwiners to a specified pure target}\label{endpoint-purity_review--classification-of-all-intertwiners-to-a-specified-pure-target}

Fix ENDPOINTSOURCEF8C6D3817D775A0F5921DFECSLOT136END, ENDPOINTSOURCEF8C6D3817D775A0F5921DFECSLOT137END, and a finite-dimensional complex vector space ENDPOINTSOURCEF8C6D3817D775A0F5921DFECSLOT138END with an invertible operator ENDPOINTSOURCEF8C6D3817D775A0F5921DFECSLOT139END. Let the source be the \textbf{original} cyclic algebra ENDPOINTSOURCEF8C6D3817D775A0F5921DFECSLOT140END, ENDPOINTSOURCEF8C6D3817D775A0F5921DFECSLOT141END, and ENDPOINTSOURCEF8C6D3817D775A0F5921DFECSLOT142END. We classify

ENDPOINTSOURCEF8C6D3817D775A0F5921DFECSLOT143END

Write ENDPOINTSOURCEF8C6D3817D775A0F5921DFECSLOT144END. The Chinese-remainder map, with the full powers, gives

ENDPOINTSOURCEF8C6D3817D775A0F5921DFECSLOT145END

Its inverse exists by Bezout: for each factor, the product of the other factors is a unit modulo that factor; multiply its inverse there and sum the resulting mutually orthogonal idempotent lifts. This retains every nilpotent power.

Let ENDPOINTSOURCEF8C6D3817D775A0F5921DFECSLOT146END be a full primary summand of ENDPOINTSOURCEF8C6D3817D775A0F5921DFECSLOT147END, with ENDPOINTSOURCEF8C6D3817D775A0F5921DFECSLOT148END any exponent at least its longest Jordan length. When ENDPOINTSOURCEF8C6D3817D775A0F5921DFECSLOT149END, define the nilpotent polynomial

ENDPOINTSOURCEF8C6D3817D775A0F5921DFECSLOT150END

The expression is independent of enlarging ENDPOINTSOURCEF8C6D3817D775A0F5921DFECSLOT151END, because additional powers vanish. The formal identities ENDPOINTSOURCEF8C6D3817D775A0F5921DFECSLOT152END and ENDPOINTSOURCEF8C6D3817D775A0F5921DFECSLOT153END, valid in the formal power-series ring, become finite polynomial identities after substituting a nilpotent. Therefore

ENDPOINTSOURCEF8C6D3817D775A0F5921DFECSLOT154END

The source satisfies ENDPOINTSOURCEF8C6D3817D775A0F5921DFECSLOT155END on ENDPOINTSOURCEF8C6D3817D775A0F5921DFECSLOT156END. Intertwining carries this equation to ENDPOINTSOURCEF8C6D3817D775A0F5921DFECSLOT157END. Projection to any different target primary summand gives zero because ENDPOINTSOURCEF8C6D3817D775A0F5921DFECSLOT158END is invertible there. Hence ENDPOINTSOURCEF8C6D3817D775A0F5921DFECSLOT159END, interpreted as zero if that primary summand is absent. Applying the logarithm polynomial to the intertwining equation proves

ENDPOINTSOURCEF8C6D3817D775A0F5921DFECSLOT160END

Conversely DP.17 implies the exponential intertwining by DP.16. Since ENDPOINTSOURCEF8C6D3817D775A0F5921DFECSLOT161END generates the source block under ENDPOINTSOURCEF8C6D3817D775A0F5921DFECSLOT162END, every such map is determined by

ENDPOINTSOURCEF8C6D3817D775A0F5921DFECSLOT163END

These are inverse constructions. Consequently

ENDPOINTSOURCEF8C6D3817D775A0F5921DFECSLOT164END

No injectivity of ENDPOINTSOURCEF8C6D3817D775A0F5921DFECSLOT165END is assumed. If two roots differ by ENDPOINTSOURCEF8C6D3817D775A0F5921DFECSLOT166END, their source blocks have the same exponential eigenvalue and map independently into the same target primary space. Their images can overlap. Formula DP.18 retains those overlaps and therefore does not incorrectly identify the total kernel with the sum of the individual block kernels for an arbitrary fixed ENDPOINTSOURCEF8C6D3817D775A0F5921DFECSLOT167END.

If ENDPOINTSOURCEF8C6D3817D775A0F5921DFECSLOT168END is pure of weight ENDPOINTSOURCEF8C6D3817D775A0F5921DFECSLOT169END in its specified complex realization, every ENDPOINTSOURCEF8C6D3817D775A0F5921DFECSLOT170END has modulus ENDPOINTSOURCEF8C6D3817D775A0F5921DFECSLOT171END. For ENDPOINTSOURCEF8C6D3817D775A0F5921DFECSLOT172END, ENDPOINTSOURCEF8C6D3817D775A0F5921DFECSLOT173END differs, so that target primary space is absent. It follows that

ENDPOINTSOURCEF8C6D3817D775A0F5921DFECSLOT174END

This proves the delivered note's off-critical annihilation including every generalized vector. Purity alone does not assert semisimplicity of ENDPOINTSOURCEF8C6D3817D775A0F5921DFECSLOT175END; DP.15--DP.19 therefore retain all target nilpotents as well.

There is also an exact quotient visible to \textbf{all} such maps. Let ENDPOINTSOURCEF8C6D3817D775A0F5921DFECSLOT176END be the largest Jordan length of ENDPOINTSOURCEF8C6D3817D775A0F5921DFECSLOT177END on ENDPOINTSOURCEF8C6D3817D775A0F5921DFECSLOT178END, and set ENDPOINTSOURCEF8C6D3817D775A0F5921DFECSLOT179END when that summand is absent. The nilpotent ENDPOINTSOURCEF8C6D3817D775A0F5921DFECSLOT180END has the same Jordan lengths as ENDPOINTSOURCEF8C6D3817D775A0F5921DFECSLOT181END, because it is the latter times an invertible polynomial with nonzero constant coefficient ENDPOINTSOURCEF8C6D3817D775A0F5921DFECSLOT182END. Put ENDPOINTSOURCEF8C6D3817D775A0F5921DFECSLOT183END. Then

ENDPOINTSOURCEF8C6D3817D775A0F5921DFECSLOT184END

To prove this, DP.18 makes every map kill ENDPOINTSOURCEF8C6D3817D775A0F5921DFECSLOT185END. Conversely choose a target Jordan chain of length ENDPOINTSOURCEF8C6D3817D775A0F5921DFECSLOT186END. If ENDPOINTSOURCEF8C6D3817D775A0F5921DFECSLOT187END, choose ENDPOINTSOURCEF8C6D3817D775A0F5921DFECSLOT188END of nilpotent length exactly ENDPOINTSOURCEF8C6D3817D775A0F5921DFECSLOT189END within that chain; if ENDPOINTSOURCEF8C6D3817D775A0F5921DFECSLOT190END, take a top vector of length ENDPOINTSOURCEF8C6D3817D775A0F5921DFECSLOT191END. The vectors ENDPOINTSOURCEF8C6D3817D775A0F5921DFECSLOT192END are independent, so the map defined by this ENDPOINTSOURCEF8C6D3817D775A0F5921DFECSLOT193END has kernel exactly the displayed ideal in that source block. Set the other source-block maps to zero. This also proves DP.21 in the presence of exponential collisions. If ENDPOINTSOURCEF8C6D3817D775A0F5921DFECSLOT194END, the displayed ideal is the entire source block.

The quotient has a canonical faithful realization in a pure target without choosing individual maps:

ENDPOINTSOURCEF8C6D3817D775A0F5921DFECSLOT195END

The operator on its target is ENDPOINTSOURCEF8C6D3817D775A0F5921DFECSLOT196END. For every ENDPOINTSOURCEF8C6D3817D775A0F5921DFECSLOT197END, ENDPOINTSOURCEF8C6D3817D775A0F5921DFECSLOT198END, hence ENDPOINTSOURCEF8C6D3817D775A0F5921DFECSLOT199END. The target is a direct sum of ENDPOINTSOURCEF8C6D3817D775A0F5921DFECSLOT200END copies of ENDPOINTSOURCEF8C6D3817D775A0F5921DFECSLOT201END, so it is pure of the same weight. Thus DP.22 induces an injection of ENDPOINTSOURCEF8C6D3817D775A0F5921DFECSLOT202END onto its actual image. This construction is the maximal quotient detected by all maps of the required type to the specified ENDPOINTSOURCEF8C6D3817D775A0F5921DFECSLOT203END. It supplies a complete typed relation rather than replacing a missing faithful map by an assertion of unrelatedness.

In particular, even a critical block can be invisible when its exponential eigenvalue is absent from ENDPOINTSOURCEF8C6D3817D775A0F5921DFECSLOT204END. A semisimple target has ENDPOINTSOURCEF8C6D3817D775A0F5921DFECSLOT205END, so all higher critical nilpotents are killed by DP.21. A nonsemisimple pure target can preserve them up to its actual Jordan length. No reduction to simple zeros is justified by purity alone.

\subsubsection{5. A weight-matched tensor construction from the original arithmetic packet}\label{endpoint-purity_review--a-weight-matched-tensor-construction-from-the-original-arithmetic-packet}

Let ENDPOINTSOURCEF8C6D3817D775A0F5921DFECSLOT206END be the original finite packet polynomial, with its full actual zero orders, and retain

ENDPOINTSOURCEF8C6D3817D775A0F5921DFECSLOT207END

The original packet uses the product of its monic linear powers, so ENDPOINTSOURCEF8C6D3817D775A0F5921DFECSLOT208END is monic in this construction. If an earlier presentation carried a nonzero scalar leading factor, the ideal map is multiplication by that recorded scalar and the factor in ENDPOINTSOURCEF8C6D3817D775A0F5921DFECSLOT209END must be retained; no such factor is set to one here. Complete zero orders make ENDPOINTSOURCEF8C6D3817D775A0F5921DFECSLOT210END a unit in the actual finite algebra. Its inverse is ENDPOINTSOURCEF8C6D3817D775A0F5921DFECSLOT211END, with the quotient analytically filled at every selected zero.

For the fixed tensor degree ENDPOINTSOURCEF8C6D3817D775A0F5921DFECSLOT212END, let

ENDPOINTSOURCEF8C6D3817D775A0F5921DFECSLOT213END

Let ENDPOINTSOURCEF8C6D3817D775A0F5921DFECSLOT214END be the original minimal annihilator of the cyclic vector ENDPOINTSOURCEF8C6D3817D775A0F5921DFECSLOT215END under ENDPOINTSOURCEF8C6D3817D775A0F5921DFECSLOT216END. This is the sum polynomial already used in the source, with all its multiplicities. Polynomial evaluation gives the injective map

ENDPOINTSOURCEF8C6D3817D775A0F5921DFECSLOT217END

Indeed its kernel in ENDPOINTSOURCEF8C6D3817D775A0F5921DFECSLOT218END is the ideal generated by that minimal annihilator, proving injectivity on the displayed quotient. Multiplication by ENDPOINTSOURCEF8C6D3817D775A0F5921DFECSLOT219END gives ENDPOINTSOURCEF8C6D3817D775A0F5921DFECSLOT220END, where ENDPOINTSOURCEF8C6D3817D775A0F5921DFECSLOT221END. Let ENDPOINTSOURCEF8C6D3817D775A0F5921DFECSLOT222END. Each tensor factor of this unit is a polynomial in the corresponding ENDPOINTSOURCEF8C6D3817D775A0F5921DFECSLOT223END; hence ENDPOINTSOURCEF8C6D3817D775A0F5921DFECSLOT224END commutes with ENDPOINTSOURCEF8C6D3817D775A0F5921DFECSLOT225END. It is invertible with inverse ENDPOINTSOURCEF8C6D3817D775A0F5921DFECSLOT226END. Therefore

ENDPOINTSOURCEF8C6D3817D775A0F5921DFECSLOT227END

This is precisely the original arithmetic inclusion. It retains the full Taylor unit rather than replacing it by a basis convention. At the source level the already defined original map satisfies

ENDPOINTSOURCEF8C6D3817D775A0F5921DFECSLOT228END

In particular the arithmetic source being embedded is the original one.

Apply DP.1--DP.4 to the original degree-ENDPOINTSOURCEF8C6D3817D775A0F5921DFECSLOT229END polynomial ENDPOINTSOURCEF8C6D3817D775A0F5921DFECSLOT230END, with independent parameters ENDPOINTSOURCEF8C6D3817D775A0F5921DFECSLOT231END and the same nonzero parameter ENDPOINTSOURCEF8C6D3817D775A0F5921DFECSLOT232END. Tensor the resulting coefficient complexes over ENDPOINTSOURCEF8C6D3817D775A0F5921DFECSLOT233END. The total differential has the usual signs: on a homogeneous tensor it differentiates the ENDPOINTSOURCEF8C6D3817D775A0F5921DFECSLOT234END-th factor with sign ENDPOINTSOURCEF8C6D3817D775A0F5921DFECSLOT235END, where ENDPOINTSOURCEF8C6D3817D775A0F5921DFECSLOT236END are the preceding cochain degrees. The direct decomposition in DP.3 contracts each factor onto its free degree-one remainder module: its other summand is an isomorphism from its degree-zero module onto its own differential image. Tensoring these decompositions leaves exactly the tensor of the degree-one remainders in degree ENDPOINTSOURCEF8C6D3817D775A0F5921DFECSLOT237END; every summand containing one of the isomorphism complexes is contracted using its inverse, with the same total-complex signs. Thus the tensor cohomology is free of rank ENDPOINTSOURCEF8C6D3817D775A0F5921DFECSLOT238END, concentrated in degree ENDPOINTSOURCEF8C6D3817D775A0F5921DFECSLOT239END, and its special fibre is ENDPOINTSOURCEF8C6D3817D775A0F5921DFECSLOT240END.

The parameter connections commute on distinct factors. Their recovered sum is

ENDPOINTSOURCEF8C6D3817D775A0F5921DFECSLOT241END

Here the bracket means insertion in that tensor factor. This is the sum, with no division by ENDPOINTSOURCEF8C6D3817D775A0F5921DFECSLOT242END.

Let ENDPOINTSOURCEF8C6D3817D775A0F5921DFECSLOT243END be the full period matrix in the original ray ordering, including every unsquared phase. Fix ENDPOINTSOURCEF8C6D3817D775A0F5921DFECSLOT244END in a sector with the stated choice of ENDPOINTSOURCEF8C6D3817D775A0F5921DFECSLOT245END; the labelled rays have angles ENDPOINTSOURCEF8C6D3817D775A0F5921DFECSLOT246END, and ENDPOINTSOURCEF8C6D3817D775A0F5921DFECSLOT247END, ENDPOINTSOURCEF8C6D3817D775A0F5921DFECSLOT248END. Each ENDPOINTSOURCEF8C6D3817D775A0F5921DFECSLOT249END is an arbitrary complex parameter. Its invertibility follows directly as follows. At the monomial potential the ENDPOINTSOURCEF8C6D3817D775A0F5921DFECSLOT250END period is a nonzero Gamma/ray factor times ENDPOINTSOURCEF8C6D3817D775A0F5921DFECSLOT251END, with ENDPOINTSOURCEF8C6D3817D775A0F5921DFECSLOT252END. A column dependence would give a polynomial of degree at most ENDPOINTSOURCEF8C6D3817D775A0F5921DFECSLOT253END vanishing at all ENDPOINTSOURCEF8C6D3817D775A0F5921DFECSLOT254END roots of unity, hence every coefficient is zero. Along a straight coefficient path the periods satisfy ENDPOINTSOURCEF8C6D3817D775A0F5921DFECSLOT255END, because integration of the image of DP.2 is an exact derivative and the endpoints vanish. Uniform domination by a negative leading exponential times a polynomial justifies each derivative. The adjugate determinant identity yields ENDPOINTSOURCEF8C6D3817D775A0F5921DFECSLOT256END, so its initially nonzero determinant never vanishes. This argument includes repeated roots.

Consequently the map

ENDPOINTSOURCEF8C6D3817D775A0F5921DFECSLOT257END

is invertible, with inverse the tensor of the exact inverses. The original arithmetic cyclic source therefore has the explicit injective rectangular period map

ENDPOINTSOURCEF8C6D3817D775A0F5921DFECSLOT258END

Holding ENDPOINTSOURCEF8C6D3817D775A0F5921DFECSLOT259END fixed in its original marked coordinates, the individual period equations and DP.26 give the exact first-order identity at ENDPOINTSOURCEF8C6D3817D775A0F5921DFECSLOT260END:

ENDPOINTSOURCEF8C6D3817D775A0F5921DFECSLOT261END

The nilpotent blocks are retained because DP.30 is injective and DP.31 intertwines the displayed recovered action. No invariance of its image under all deformed ENDPOINTSOURCEF8C6D3817D775A0F5921DFECSLOT262END, and no missing higher-order period identity, is assumed.

For any original finite-cutoff positive Gram ENDPOINTSOURCEF8C6D3817D775A0F5921DFECSLOT263END, put ENDPOINTSOURCEF8C6D3817D775A0F5921DFECSLOT264END and ENDPOINTSOURCEF8C6D3817D775A0F5921DFECSLOT265END. Then ENDPOINTSOURCEF8C6D3817D775A0F5921DFECSLOT266END, and the concrete positive form on ENDPOINTSOURCEF8C6D3817D775A0F5921DFECSLOT267END is the restriction of ENDPOINTSOURCEF8C6D3817D775A0F5921DFECSLOT268END. Its pullback is

ENDPOINTSOURCEF8C6D3817D775A0F5921DFECSLOT269END

The action on that image is ENDPOINTSOURCEF8C6D3817D775A0F5921DFECSLOT270END, and its control form pulls back to ENDPOINTSOURCEF8C6D3817D775A0F5921DFECSLOT271END. Thus every original mass, source norm, and generalized defect is preserved by an actual metric-and-action map. The ambient form may be degenerate outside the image; no positive extension or orthogonal-complement invariance is asserted.

There is a finite coefficient model retaining this injection. Form the finitely generated subring of ENDPOINTSOURCEF8C6D3817D775A0F5921DFECSLOT272END containing the coefficients of ENDPOINTSOURCEF8C6D3817D775A0F5921DFECSLOT273END, ENDPOINTSOURCEF8C6D3817D775A0F5921DFECSLOT274END, ENDPOINTSOURCEF8C6D3817D775A0F5921DFECSLOT275END, its inverse, and the matrices in DP.25--DP.26, together with all denominators ENDPOINTSOURCEF8C6D3817D775A0F5921DFECSLOT276END. Since ENDPOINTSOURCEF8C6D3817D775A0F5921DFECSLOT277END has full column rank, choose one of its nonzero square maximal minors ENDPOINTSOURCEF8C6D3817D775A0F5921DFECSLOT278END and adjoin ENDPOINTSOURCEF8C6D3817D775A0F5921DFECSLOT279END. Call the resulting ring ENDPOINTSOURCEF8C6D3817D775A0F5921DFECSLOT280END. The selected rows give a left inverse by their adjugate divided by ENDPOINTSOURCEF8C6D3817D775A0F5921DFECSLOT281END. Thus ENDPOINTSOURCEF8C6D3817D775A0F5921DFECSLOT282END is a split injection over this very ring, and its action identity holds there because every entry is an element of its subring in ENDPOINTSOURCEF8C6D3817D775A0F5921DFECSLOT283END and the equality holds in ENDPOINTSOURCEF8C6D3817D775A0F5921DFECSLOT284END.

For each specified reduction ENDPOINTSOURCEF8C6D3817D775A0F5921DFECSLOT285END with ENDPOINTSOURCEF8C6D3817D775A0F5921DFECSLOT286END, the same matrices remain a split injection and intertwiner of the reduced coefficient algebras. Its reduced cyclic image is therefore isomorphic, through those matrices, to the literal quotient ENDPOINTSOURCEF8C6D3817D775A0F5921DFECSLOT287END. Every nilpotent power of that reduced image is the corresponding power in this quotient, including any root collisions after reduction. The nonzero minor used for this particular finite construction has been inverted. This does not claim that every prime is allowed or that one localization works for unbounded ENDPOINTSOURCEF8C6D3817D775A0F5921DFECSLOT288END.

On the finite-field branch, form the \textbf{external product} of the ENDPOINTSOURCEF8C6D3817D775A0F5921DFECSLOT289END original one-variable nontrivial additive-character sheaves. The Künneth map for compact support identifies its degree-ENDPOINTSOURCEF8C6D3817D775A0F5921DFECSLOT290END cohomology, with product-coordinate orientation, with

ENDPOINTSOURCEF8C6D3817D775A0F5921DFECSLOT291END

Each factor is concentrated in degree one and pure of weight one by the checked theorem. Thus DP.33 is concentrated in degree ENDPOINTSOURCEF8C6D3817D775A0F5921DFECSLOT292END, has rank ENDPOINTSOURCEF8C6D3817D775A0F5921DFECSLOT293END, and is pure of weight ENDPOINTSOURCEF8C6D3817D775A0F5921DFECSLOT294END: its Frobenius is the tensor product, whose eigenvalues are all products of the factor eigenvalues. To see the last assertion including nonsemisimple factors, triangularize each complex coefficient realization; the tensor basis makes the tensor product triangular, with precisely those products on its diagonal. Each product has modulus ENDPOINTSOURCEF8C6D3817D775A0F5921DFECSLOT295END. The alternating trace sign is ENDPOINTSOURCEF8C6D3817D775A0F5921DFECSLOT296END, agreeing with the product of the ENDPOINTSOURCEF8C6D3817D775A0F5921DFECSLOT297END one-variable negative trace signs.

The retained construction therefore has a weight-matched common coefficient span, with DP.28--DP.32 on its complex side and DP.33 on its finite-field side. The reduced finite coefficient injection is not, by its existence, a map into the separate ENDPOINTSOURCEF8C6D3817D775A0F5921DFECSLOT298END-adic cohomology. A specified operator-compatible comparison with that target is governed exactly by DP.13--DP.22. This locates the remaining comparison inside an explicit surviving geometric construction, preserving the arithmetic input and tensor degree throughout.

\subsubsection{5A. Placement inside the original cohomology program over the privileged tau point}\label{endpoint-purity_review--a.-placement-inside-the-original-cohomology-program-over-the-privileged-tau-point}

The preceding tensor construction is now placed in the actual base program, using the complete original \texttt{Tau\_Base\_Cohomology\_2026-09-12/NOTE.md}, read in full. It does not replace that base by a finite field.

Retain the original pointed blueprint

ENDPOINTSOURCEF8C6D3817D775A0F5921DFECSLOT299END

Its additive zero relation is that the empty sum represents ENDPOINTSOURCEF8C6D3817D775A0F5921DFECSLOT300END; no Boolean equation ENDPOINTSOURCEF8C6D3817D775A0F5921DFECSLOT301END is introduced into this base. For each original nonzero ring ENDPOINTSOURCEF8C6D3817D775A0F5921DFECSLOT302END, ENDPOINTSOURCEF8C6D3817D775A0F5921DFECSLOT303END is the blueprint of ENDPOINTSOURCEF8C6D3817D775A0F5921DFECSLOT304END with its original additive relations. A prime of ENDPOINTSOURCEF8C6D3817D775A0F5921DFECSLOT305END contains ENDPOINTSOURCEF8C6D3817D775A0F5921DFECSLOT306END and does not contain the multiplicative identity. Its inverse image under ENDPOINTSOURCEF8C6D3817D775A0F5921DFECSLOT307END is therefore exactly ENDPOINTSOURCEF8C6D3817D775A0F5921DFECSLOT308END. This proves that the structural spectrum map lands at the single privileged point through the actual prime inverse-image map. The arithmetic map ENDPOINTSOURCEF8C6D3817D775A0F5921DFECSLOT309END, ENDPOINTSOURCEF8C6D3817D775A0F5921DFECSLOT310END, and the Boolean support map ENDPOINTSOURCEF8C6D3817D775A0F5921DFECSLOT311END, ENDPOINTSOURCEF8C6D3817D775A0F5921DFECSLOT312END, remain attached. Their pair is injective, since support separates ENDPOINTSOURCEF8C6D3817D775A0F5921DFECSLOT313END from every represented element and amplitude distinguishes represented elements. The map here is DB.1; the separate localization that forces ENDPOINTSOURCEF8C6D3817D775A0F5921DFECSLOT314END is not performed.

For a coefficient map ENDPOINTSOURCEF8C6D3817D775A0F5921DFECSLOT315END, its split lift satisfies

ENDPOINTSOURCEF8C6D3817D775A0F5921DFECSLOT316END

Each equality follows on ENDPOINTSOURCEF8C6D3817D775A0F5921DFECSLOT317END and on each ENDPOINTSOURCEF8C6D3817D775A0F5921DFECSLOT318END, the two exhaustive kinds of elements. Apply DB.2 to the maps from ENDPOINTSOURCEF8C6D3817D775A0F5921DFECSLOT319END in DP.26--DP.33, to the complex coefficient realization, and to every specified finite-field reduction. All these diagrams remain over the same ENDPOINTSOURCEF8C6D3817D775A0F5921DFECSLOT320END.

There is an exact categorical version for the linear fibres. For an ENDPOINTSOURCEF8C6D3817D775A0F5921DFECSLOT321END-module ENDPOINTSOURCEF8C6D3817D775A0F5921DFECSLOT322END, set ENDPOINTSOURCEF8C6D3817D775A0F5921DFECSLOT323END, with

ENDPOINTSOURCEF8C6D3817D775A0F5921DFECSLOT324END

For an ENDPOINTSOURCEF8C6D3817D775A0F5921DFECSLOT325END-linear map ENDPOINTSOURCEF8C6D3817D775A0F5921DFECSLOT326END, its lift sends ENDPOINTSOURCEF8C6D3817D775A0F5921DFECSLOT327END to ENDPOINTSOURCEF8C6D3817D775A0F5921DFECSLOT328END and ENDPOINTSOURCEF8C6D3817D775A0F5921DFECSLOT329END to ENDPOINTSOURCEF8C6D3817D775A0F5921DFECSLOT330END. This is ENDPOINTSOURCEF8C6D3817D775A0F5921DFECSLOT331END-linear by the two ordinary linearity identities and the rules for ENDPOINTSOURCEF8C6D3817D775A0F5921DFECSLOT332END. Conversely, a ENDPOINTSOURCEF8C6D3817D775A0F5921DFECSLOT333END-linear map between these objects which preserves represented support restricts to a map ENDPOINTSOURCEF8C6D3817D775A0F5921DFECSLOT334END. Its additive and scalar identities are exactly the original ENDPOINTSOURCEF8C6D3817D775A0F5921DFECSLOT335END-linear identities. Moreover its value on ENDPOINTSOURCEF8C6D3817D775A0F5921DFECSLOT336END is ENDPOINTSOURCEF8C6D3817D775A0F5921DFECSLOT337END. Thus this gives a fully faithful equivalence with the category consisting of these split module objects and support-preserving maps.

In that category the zero morphism is the supported-zero map ENDPOINTSOURCEF8C6D3817D775A0F5921DFECSLOT338END, with ENDPOINTSOURCEF8C6D3817D775A0F5921DFECSLOT339END. The zero object is ENDPOINTSOURCEF8C6D3817D775A0F5921DFECSLOT340END. These assertions follow either by the fully faithful correspondence or directly from DB.3: each zero-object arrow preserving support is forced, and composition through it gives exactly that supported-zero map. Consequently the original equation ENDPOINTSOURCEF8C6D3817D775A0F5921DFECSLOT341END is carried to the correct supported-zero composition, preserving the separate absorbing element. This supplies the precise linear-fibre category in which the original support cohomology is reconstructed. With the several masks of the source, apply it at each admitted mask and retain all inclusion transports; do not tensor the masks as if they were unrecorded vector coordinates.

Now retain the actual chart poset

ENDPOINTSOURCEF8C6D3817D775A0F5921DFECSLOT342END

and the actual coefficient sheaf

ENDPOINTSOURCEF8C6D3817D775A0F5921DFECSLOT343END

The ENDPOINTSOURCEF8C6D3817D775A0F5921DFECSLOT344END in DB.5 is the original reflection ENDPOINTSOURCEF8C6D3817D775A0F5921DFECSLOT345END; the finite jet map below is written ENDPOINTSOURCEF8C6D3817D775A0F5921DFECSLOT346END. Let ENDPOINTSOURCEF8C6D3817D775A0F5921DFECSLOT347END for ENDPOINTSOURCEF8C6D3817D775A0F5921DFECSLOT348END and zero otherwise. Evaluation gives ENDPOINTSOURCEF8C6D3817D775A0F5921DFECSLOT349END, hence these diagrams are projective. The sequence

ENDPOINTSOURCEF8C6D3817D775A0F5921DFECSLOT350END

is exact: at ENDPOINTSOURCEF8C6D3817D775A0F5921DFECSLOT351END and ENDPOINTSOURCEF8C6D3817D775A0F5921DFECSLOT352END its last map is the identity, and at ENDPOINTSOURCEF8C6D3817D775A0F5921DFECSLOT353END and ENDPOINTSOURCEF8C6D3817D775A0F5921DFECSLOT354END its kernel is exactly the displayed anti-diagonal. Each map commutes with the poset arrows because its entries are the same scalar identity at all nonzero stalks. Applying ENDPOINTSOURCEF8C6D3817D775A0F5921DFECSLOT355END proves

ENDPOINTSOURCEF8C6D3817D775A0F5921DFECSLOT356END

On DB.5 this is the original two-leg theta complex. Its degree-zero kernel consists of ENDPOINTSOURCEF8C6D3817D775A0F5921DFECSLOT357END, since ENDPOINTSOURCEF8C6D3817D775A0F5921DFECSLOT358END is injective, and its degree-one quotient is ENDPOINTSOURCEF8C6D3817D775A0F5921DFECSLOT359END. This identifies the actual arithmetic cohomology over the privileged base point. Restriction to the ENDPOINTSOURCEF8C6D3817D775A0F5921DFECSLOT360END stalk is a separate map: its two degree-zero representatives ENDPOINTSOURCEF8C6D3817D775A0F5921DFECSLOT361END and ENDPOINTSOURCEF8C6D3817D775A0F5921DFECSLOT362END differ by ENDPOINTSOURCEF8C6D3817D775A0F5921DFECSLOT363END, so the homotopy is exactly ENDPOINTSOURCEF8C6D3817D775A0F5921DFECSLOT364END. This proves its derived comparison to DB.7 without replacing global cohomology by the support stalk.

Define ENDPOINTSOURCEF8C6D3817D775A0F5921DFECSLOT365END for ENDPOINTSOURCEF8C6D3817D775A0F5921DFECSLOT366END and zero otherwise. The identity ENDPOINTSOURCEF8C6D3817D775A0F5921DFECSLOT367END and exactness of evaluation show that each is injective over ENDPOINTSOURCEF8C6D3817D775A0F5921DFECSLOT368END. Dualizing DB.6 gives

ENDPOINTSOURCEF8C6D3817D775A0F5921DFECSLOT369END

Applying ENDPOINTSOURCEF8C6D3817D775A0F5921DFECSLOT370END gives the explicit differential ENDPOINTSOURCEF8C6D3817D775A0F5921DFECSLOT371END, which is the dual of DB.7. This proves its stated adjunction to ENDPOINTSOURCEF8C6D3817D775A0F5921DFECSLOT372END. The stalk complexes at ENDPOINTSOURCEF8C6D3817D775A0F5921DFECSLOT373END are ENDPOINTSOURCEF8C6D3817D775A0F5921DFECSLOT374END and ENDPOINTSOURCEF8C6D3817D775A0F5921DFECSLOT375END, so they are acyclic; at ENDPOINTSOURCEF8C6D3817D775A0F5921DFECSLOT376END the sole term is ENDPOINTSOURCEF8C6D3817D775A0F5921DFECSLOT377END in degree ENDPOINTSOURCEF8C6D3817D775A0F5921DFECSLOT378END, and at ENDPOINTSOURCEF8C6D3817D775A0F5921DFECSLOT379END all terms vanish. Thus ENDPOINTSOURCEF8C6D3817D775A0F5921DFECSLOT380END, with the exact inherited shift.

Tensor DB.6 over ENDPOINTSOURCEF8C6D3817D775A0F5921DFECSLOT381END in the fixed factor order to resolve the constant diagram on ENDPOINTSOURCEF8C6D3817D775A0F5921DFECSLOT382END. Its differential is the signed sum with sign ENDPOINTSOURCEF8C6D3817D775A0F5921DFECSLOT383END. In top degree the actual coefficient numerator is ENDPOINTSOURCEF8C6D3817D775A0F5921DFECSLOT384END, and the boundary subspace is the sum of tensors with ENDPOINTSOURCEF8C6D3817D775A0F5921DFECSLOT385END in at least one factor. Quotienting one factor at a time proves, by the defining relations of the tensor quotient,

ENDPOINTSOURCEF8C6D3817D775A0F5921DFECSLOT386END

The second identity follows by the same dual tensor resolution: any stalk with at least one ENDPOINTSOURCEF8C6D3817D775A0F5921DFECSLOT387END or ENDPOINTSOURCEF8C6D3817D775A0F5921DFECSLOT388END factor contains its signed two-term identity contraction, any ENDPOINTSOURCEF8C6D3817D775A0F5921DFECSLOT389END factor gives zero, and the all-ENDPOINTSOURCEF8C6D3817D775A0F5921DFECSLOT390END stalk has exactly ENDPOINTSOURCEF8C6D3817D775A0F5921DFECSLOT391END in degree ENDPOINTSOURCEF8C6D3817D775A0F5921DFECSLOT392END. This proves the product right-adjoint calculation with its signs and no deleted cohomological shift.

Use the actual full-jet map ENDPOINTSOURCEF8C6D3817D775A0F5921DFECSLOT393END. Since ENDPOINTSOURCEF8C6D3817D775A0F5921DFECSLOT394END, its jets vanish through every full selected zero order; only holomorphy in neighborhoods of those selected zeros in the strip is required here. The map with component ENDPOINTSOURCEF8C6D3817D775A0F5921DFECSLOT395END at ENDPOINTSOURCEF8C6D3817D775A0F5921DFECSLOT396END and zero at the other stalks is therefore a sheaf map ENDPOINTSOURCEF8C6D3817D775A0F5921DFECSLOT397END. Its tensor product gives the actual top-cohomology arrow

ENDPOINTSOURCEF8C6D3817D775A0F5921DFECSLOT398END

It is surjective because ENDPOINTSOURCEF8C6D3817D775A0F5921DFECSLOT399END in the original full-order packet; tensoring that explicit right inverse gives a right inverse of DB.10. Compose the other way with the original cyclic inclusion to obtain the arithmetic split pair

ENDPOINTSOURCEF8C6D3817D775A0F5921DFECSLOT400END

Thus DP.30 is the period realization of the actual cyclic submodule of this tau-base cohomology: the map is ENDPOINTSOURCEF8C6D3817D775A0F5921DFECSLOT401END restricted through DB.11. The units in ENDPOINTSOURCEF8C6D3817D775A0F5921DFECSLOT402END, the full jets, and the original cochain representative boundaries are all present. On top cohomology the scaling action is ENDPOINTSOURCEF8C6D3817D775A0F5921DFECSLOT403END; its exact intertwining through DB.10 and DB.11 is ENDPOINTSOURCEF8C6D3817D775A0F5921DFECSLOT404END. This follows from DP.26 by the convergent exponential series, or by its finite nilpotent expansions on each block.

The original right-adjoint dual injection also extends to this cyclic module. For the original reflection-stable packet ENDPOINTSOURCEF8C6D3817D775A0F5921DFECSLOT405END, let ENDPOINTSOURCEF8C6D3817D775A0F5921DFECSLOT406END be the full residue form from the tau-base note,

ENDPOINTSOURCEF8C6D3817D775A0F5921DFECSLOT407END

Its local denominator is the actual ENDPOINTSOURCEF8C6D3817D775A0F5921DFECSLOT408END, ENDPOINTSOURCEF8C6D3817D775A0F5921DFECSLOT409END. Multiplication by ENDPOINTSOURCEF8C6D3817D775A0F5921DFECSLOT410END is an invertible triangular Taylor map; extracting the coefficient of ENDPOINTSOURCEF8C6D3817D775A0F5921DFECSLOT411END pairs powers ENDPOINTSOURCEF8C6D3817D775A0F5921DFECSLOT412END with nonzero anti-diagonal. Reflection sends the opposite block's ENDPOINTSOURCEF8C6D3817D775A0F5921DFECSLOT413END to ENDPOINTSOURCEF8C6D3817D775A0F5921DFECSLOT414END in coefficients. Hence this sesquilinear form is nondegenerate on the full finite direct sum, retaining all nilpotents and the exact reflection signs. Tensor its numerical values in the fixed order to get a nondegenerate sesquilinear pairing ENDPOINTSOURCEF8C6D3817D775A0F5921DFECSLOT415END on ENDPOINTSOURCEF8C6D3817D775A0F5921DFECSLOT416END; the tensor Gram is the tensor of invertible Grams and has the tensor of their inverses. No phase is inserted to claim positivity. In particular the skew-Hermitian parity is retained as ENDPOINTSOURCEF8C6D3817D775A0F5921DFECSLOT417END under conjugate exchange.

Let ENDPOINTSOURCEF8C6D3817D775A0F5921DFECSLOT418END be the original moment line, with scaling action ENDPOINTSOURCEF8C6D3817D775A0F5921DFECSLOT419END. Define

ENDPOINTSOURCEF8C6D3817D775A0F5921DFECSLOT420END

The displayed cohomological identification follows from the adjunction DB.8--DB.9: over the field, taking the algebraic dual is exact, so degree ENDPOINTSOURCEF8C6D3817D775A0F5921DFECSLOT421END is the dual of top degree ENDPOINTSOURCEF8C6D3817D775A0F5921DFECSLOT422END. The formula is linear in ENDPOINTSOURCEF8C6D3817D775A0F5921DFECSLOT423END. It is well-defined on ENDPOINTSOURCEF8C6D3817D775A0F5921DFECSLOT424END because DB.10 is, including every theta relation. If it vanishes for all ENDPOINTSOURCEF8C6D3817D775A0F5921DFECSLOT425END, surjectivity of DB.10 and nondegeneracy of the tensor residue form imply ENDPOINTSOURCEF8C6D3817D775A0F5921DFECSLOT426END; injectivity of DP.26 then implies ENDPOINTSOURCEF8C6D3817D775A0F5921DFECSLOT427END. Thus DB.13 is an actual injection to the computed tau-base dual object.

For real ENDPOINTSOURCEF8C6D3817D775A0F5921DFECSLOT428END, the original scaling multiplier satisfies ENDPOINTSOURCEF8C6D3817D775A0F5921DFECSLOT429END. Multiplying the two factors in DB.12 therefore gives exactly ENDPOINTSOURCEF8C6D3817D775A0F5921DFECSLOT430END, so ENDPOINTSOURCEF8C6D3817D775A0F5921DFECSLOT431END, including the full Taylor expansions. Tensoring gives the factor ENDPOINTSOURCEF8C6D3817D775A0F5921DFECSLOT432END. On the algebraic dual the action is ENDPOINTSOURCEF8C6D3817D775A0F5921DFECSLOT433END. It follows by direct substitution in DB.13 that this injection is equivariant with the conjugate of the original cyclic scaling action. This proves the tensor counterpart of the original tau-base note's (43), with ENDPOINTSOURCEF8C6D3817D775A0F5921DFECSLOT434END and the correct moment line, rather than leaving the pairing outside the cohomological model.

Finally, each nonempty original support mask retains its copy of the coefficient complex and each transport retains its stated linear map. Reconstructing them by DB.3 gives the supported degree-one object ENDPOINTSOURCEF8C6D3817D775A0F5921DFECSLOT435END and the original tensor supported cohomology. All maps DB.10--DB.13 preserve the admitted mask; a zero amplitude maps to that mask's supported zero and external absence maps to ENDPOINTSOURCEF8C6D3817D775A0F5921DFECSLOT436END. The dual mask arrows are precomposition in the opposite mask category, as in the original note. This completes the placement of the weight-matched auxiliary tensor comparison \textbf{inside} the actual split cohomology program over ENDPOINTSOURCEF8C6D3817D775A0F5921DFECSLOT437END.

\subsubsection{6. Original critical-line comparison and the retained kernel}\label{endpoint-purity_review--original-critical-line-comparison-and-the-retained-kernel}

The complete local proof of the critical-only comparison and of the arbitrary-polynomial theta sequence is supplied in the companion \texttt{deligne\_bound\_critical\_kernel\_subreview\_20260913.md}. Its sources were read in full independently. The exact sequence from the original theta complex is

ENDPOINTSOURCEF8C6D3817D775A0F5921DFECSLOT438END

The first arrow sends a class represented by ENDPOINTSOURCEF8C6D3817D775A0F5921DFECSLOT439END, with ENDPOINTSOURCEF8C6D3817D775A0F5921DFECSLOT440END, to ENDPOINTSOURCEF8C6D3817D775A0F5921DFECSLOT441END. It is independent of the representative because changing ENDPOINTSOURCEF8C6D3817D775A0F5921DFECSLOT442END by ENDPOINTSOURCEF8C6D3817D775A0F5921DFECSLOT443END changes ENDPOINTSOURCEF8C6D3817D775A0F5921DFECSLOT444END by ENDPOINTSOURCEF8C6D3817D775A0F5921DFECSLOT445END. The source and target Euler inverse formulas in \texttt{arithmetic\_input.tex} show that this arrow identifies ENDPOINTSOURCEF8C6D3817D775A0F5921DFECSLOT446END with the kernel of the displayed multiplication. The target jet map identifies its cokernel with ENDPOINTSOURCEF8C6D3817D775A0F5921DFECSLOT447END. This retains the exact theta source, not an arbitrary test spectral model.

At a selected centre, write ENDPOINTSOURCEF8C6D3817D775A0F5921DFECSLOT448END, ENDPOINTSOURCEF8C6D3817D775A0F5921DFECSLOT449END, ENDPOINTSOURCEF8C6D3817D775A0F5921DFECSLOT450END, with ENDPOINTSOURCEF8C6D3817D775A0F5921DFECSLOT451END, and ENDPOINTSOURCEF8C6D3817D775A0F5921DFECSLOT452END; when ENDPOINTSOURCEF8C6D3817D775A0F5921DFECSLOT453END, set ENDPOINTSOURCEF8C6D3817D775A0F5921DFECSLOT454END. In ENDPOINTSOURCEF8C6D3817D775A0F5921DFECSLOT455END, multiplication by ENDPOINTSOURCEF8C6D3817D775A0F5921DFECSLOT456END has kernel ENDPOINTSOURCEF8C6D3817D775A0F5921DFECSLOT457END and cokernel ENDPOINTSOURCEF8C6D3817D775A0F5921DFECSLOT458END. Indeed the factor ENDPOINTSOURCEF8C6D3817D775A0F5921DFECSLOT459END is invertible and multiplication by ENDPOINTSOURCEF8C6D3817D775A0F5921DFECSLOT460END has precisely these kernel and image ideals. These formulas prove that vanishing occurs whenever ENDPOINTSOURCEF8C6D3817D775A0F5921DFECSLOT461END, including undercounted orders. Exact full order ENDPOINTSOURCEF8C6D3817D775A0F5921DFECSLOT462END makes ENDPOINTSOURCEF8C6D3817D775A0F5921DFECSLOT463END a unit; it is that unit which appears in DP.23 and the full packet retraction. The precise local map from the kernel to the quotient, through the original cohomology, is proved in the companion, so that no equality between two different jet presentations is presumed here.

For a full actual packet ENDPOINTSOURCEF8C6D3817D775A0F5921DFECSLOT464END, the existing Connes comparison has

ENDPOINTSOURCEF8C6D3817D775A0F5921DFECSLOT465END

The full higher-jet inverse on the selected critical blocks identifies its actual image with ENDPOINTSOURCEF8C6D3817D775A0F5921DFECSLOT466END. Thus DP.35 yields an exact equivariant sequence

ENDPOINTSOURCEF8C6D3817D775A0F5921DFECSLOT467END

with the canonical primary splitting of its source. Its two-term comparison complex ENDPOINTSOURCEF8C6D3817D775A0F5921DFECSLOT468END, in degrees zero and one, has degree-zero cohomology ENDPOINTSOURCEF8C6D3817D775A0F5921DFECSLOT469END and degree-one cohomology zero. The inclusion of the kernel in degree zero is therefore a quasi-isomorphism. If the full target is retained instead, degree-one cohomology is its quotient by ENDPOINTSOURCEF8C6D3817D775A0F5921DFECSLOT470END; it must not be silently set to zero. This is an exact derived record of everything killed by the critical comparison.

For any receiving vector-space fibre labelled by ENDPOINTSOURCEF8C6D3817D775A0F5921DFECSLOT471END, the split lift of a linear map is

ENDPOINTSOURCEF8C6D3817D775A0F5921DFECSLOT472END

The preimages are exactly

ENDPOINTSOURCEF8C6D3817D775A0F5921DFECSLOT473END

Consequently DP.35 retains its full off-critical kernel as supported source vectors. The support observation distinguishes those vectors from external absence; it does not turn an annihilating linear map into an injection. At the same time DP.36 gives the precise surviving morphism and the complex that records the lost blocks. That is a structural continuation of the program, not a refutation of it.

\subsubsection{7. Acceptance and the next exact arithmetic calculation}\label{endpoint-purity_review--acceptance-and-the-next-exact-arithmetic-calculation}

Equations (12)--(13), (17)--(24), (26)--(27), and (44)--(50) of the delivered note have the stated mathematical scope after the character-summand clarification. The note's full-order sentence following (57) needs the distinction proved in DP.34. Its quantitative Gamma estimate is covered by the separate independent Gamma review; this document does not convert a full reading of that section into an independent numerical or all-ENDPOINTSOURCEF8C6D3817D775A0F5921DFECSLOT474END certificate.

The note's period congruence is exact: with a common invertible ENDPOINTSOURCEF8C6D3817D775A0F5921DFECSLOT475END, ENDPOINTSOURCEF8C6D3817D775A0F5921DFECSLOT476END and ENDPOINTSOURCEF8C6D3817D775A0F5921DFECSLOT477END give ENDPOINTSOURCEF8C6D3817D775A0F5921DFECSLOT478END by direct multiplication. In particular every generalized eigenvalue and every common-index determinant ratio is retained. This is consistent with the faithful tensor map DP.30 and the original theta relation DP.27.

The actual arithmetic multiplier still enters the two source metrics through their identity coefficient map and their different least-norm sections. Its signed endpoint determinant correction in delivered (54) remains an exact quantity in the program. The present audit neither assumes that correction is small nor replaces its calculation by a condition. It establishes the full comparison diagram that remains valid while that source-specific calculation proceeds.

The phrase ``the F1 geometry program has failed'' does not follow from this package. A defensible conclusion is: the displayed auxiliary purity, by itself and through the existing maps, does not force the desired original theta-source estimate. The exact failures and surviving maps have now been calculated, including a weight-ENDPOINTSOURCEF8C6D3817D775A0F5921DFECSLOT479END tensor construction and the complete target-dependent comparison kernel.

\subsubsection{Source pins and actual reading record}\label{endpoint-purity_review--source-pins-and-actual-reading-record}

The companion JSON records SHA-256, byte counts, and reading scope for every selected source, this review, and the independent critical-kernel proof. The new delivered note was read completely. \texttt{arithmetic\_input.tex}, \texttt{Exponential\_Translation\_Comparison.tex}, and the primary correction were read completely. The SGA trace-period note was read from Sections 6 through 12, including all coefficient, period, tensor-weight, and metric formulas used above. The original Deligne scan was visually inspected on complete printed pages 215--216 only. The retained S20 text is a location aid; the original scan supplies the theorem typography.

No Lean invocation, remote mutation, whole-library audit, or numerical certification is claimed in this review.
